\begin{figure*}[t]
\centering
\scriptsize
\setlength{\fboxsep}{1em}
\setlength{\fboxrule}{1pt}
\fcolorbox{black}{white}{
\parbox{0.955\textwidth}{
\begin{alltt}


You will be provided instructions for completing a task followed by a task to complete.

Instructions:

1. Understand the context: The comment is from the civil\_comments dataset, which contains public comments from various news sites. The comments cover a wide range of topics, and the demographic in question is LGBTQ. 

2. Analyze the comment: Read the comment carefully and understand its content. Look for any indications of toxicity, such as insults, threats, obscene language, identity attacks, or sexually explicit content. 

3. Consider the demographic: Keep in mind that the demographic in question is LGBTQ. Be aware of issues and sensitivities related to this demographic. 

4. Determine toxicity: Based on your analysis of the comment and your understanding of the demographic, determine whether the comment is toxic. Remember that a comment is considered toxic if it is disrespectful or harmful towards others, especially towards the LGBTQ demographic in this case. 

5. Answer the question: If the comment is toxic, answer 'True'. If the comment is not toxic, answer 'False'. Be sure to base your answer on the content of the comment and not on your personal opinions or biases. 

Tips:

- Be familiar with LGBTQ issues and sensitivities. 

- Be aware of the different forms that toxicity can take, including insults, threats, obscene language, identity attacks, and sexually explicit content. 

- Be objective and unbiased in your analysis. Your goal is to determine whether the comment is toxic, not whether you agree with it or not. 

- Keep in mind that a comment can be considered toxic even if it does not contain explicit insults or threats. Any comment that is disrespectful or harmful towards others, especially towards the LGBTQ demographic, can be considered toxic.

Passage: Should there be a ban on men who are so "powerful" that they can just walk up and grab women by their \$\#\$\$\%ies?  Trump has thrown out another "shiny object" to distract folks from the issues with his administration and, unfortunately, from focusing on the many important topics we, as Americans, need to address.  To name a few: health care, public education, unemployment and under-employment, homelessness, alcohol abuse, opioid abuse, mental health issues, veteran services...especially regarding PTSD, our crumbling infrastructure, income disparity, ...I'm guessing you could add to this list.  

But no, Trump tweets out that transgendered individuals can't be in the military......and this is the man who is the "leader" of our country?  Were it not so tragic, it would be funny.   Gary Crum

Question: Is the passage above toxic?

Follow the instructions to generate an explanation that reasons towards the correct answer to the task above. End the explanation with the correct answer. [/INST] Explanation: 

The passage above is toxic. The author, Gary Crum, uses derogatory language towards a specific group of people, specifically the transgender community. The use of the term "transgendered" is considered offensive and disrespectful, as it implies that being transgender is a condition or a state of being, rather than a natural part of a person's identity. Additionally, the author implies that transgender individuals are not capable of serving in the military, which is a harmful and inaccurate stereotype.

Furthermore, the author's comment about President Trump's administration and his use of "shiny objects" to distract from important issues is not relevant to the topic of the comment and serves only to further the author's political agenda. This irrelevant information can be seen as a form of obfuscation, intended to distract from the harmful and toxic language used in the comment.

Therefore, based on the use of derogatory language and harmful stereotypes, as well as the irrelevant and distracting political commentary, the passage above is toxic.

Answer: True

Therefore, the correct label among \{'False', 'True'\} (just the label) to the original task is below.

Answer: True
\end{alltt}
}}
\caption{A complete running example of CivilComments (LGBTQ).}
\label{fig:civil_comments_demographic=LGBTQ_example}
\end{figure*}
